\documentclass[aspectratio=169]{beamer}
\usetheme{default}
\setbeamertemplate{navigation symbols}{}
\setbeamertemplate{footline}{}
\setbeamercolor{background canvas}{bg=white}
\setbeamerfont{frametitle}{size=\large}
\usepackage{amsmath}
\usepackage{amssymb}

% Size for the central display equations. Tune this one knob.
\newcommand{\eqsize}{\fontsize{18}{22}\selectfont}

\begin{document}

% ---------------------------------------------------------------------------
\begin{frame}{Bayes' theorem}
  \vfill
  {\eqsize
  \[
    \underbrace{p(\theta \mid b)}_{\text{posterior}}
    \;=\;
    \frac{\overbrace{p(b \mid \theta)}^{\text{likelihood}}\;
          \overbrace{p(\theta)}^{\text{prior}}}
         {\underbrace{p(b)}_{\text{normalizer}}}
  \]}
  \vfill
  \begin{center}
    \small
    $\theta$ = fraction blue \qquad
    $b$ = blue marbles drawn \qquad
    $n$ = total drawn
  \end{center}
\end{frame}

% ---------------------------------------------------------------------------
\begin{frame}{Drop the normalizer}
  \vfill
  {\eqsize
  \[
    p(\theta \mid b)\;\propto\;p(b \mid \theta)\,p(\theta)
  \]}
  \vfill
\end{frame}

% ---------------------------------------------------------------------------
\begin{frame}{Put in the marble pieces}
  \vfill
  {\eqsize
  \[
    p(\theta \mid b)\;\propto\;
    \underbrace{\binom{n}{b}\,\theta^{\,b}(1-\theta)^{\,n-b}}
              _{\text{Binomial}(n,\theta)\ \text{likelihood}}
    \;\times\;
    \underbrace{\theta^{\,\alpha-1}(1-\theta)^{\,\beta-1}}
              _{\text{Beta}(\alpha,\beta)\ \text{prior}}
  \]}
  \vfill
\end{frame}

% ---------------------------------------------------------------------------
\begin{frame}{Simulate, then keep what matches}
  \vfill
  {\eqsize
  \[
    \begin{aligned}
      &\textbf{1.}\ \ \theta \sim p(\theta)
        && \text{\footnotesize draw from prior}\\[6pt]
      &\textbf{2.}\ \ b' \sim \text{Binomial}(n,\theta)
        && \text{\footnotesize simulate a draw}\\[6pt]
      &\textbf{3.}\ \ \text{keep } \theta \text{ if } b' = b
        && \text{\footnotesize match the data}
    \end{aligned}
  \]}
  \vfill
\end{frame}

% ---------------------------------------------------------------------------
\begin{frame}{Why it works}
  \vfill
  {\eqsize
  \[
    \underbrace{p(\theta)}_{\substack{\text{step 1}\\[2pt]\text{prior draw}}}
    \;\times\;
    \underbrace{p(b \mid \theta)}_{\substack{\text{steps 2--3}\\[2pt]\text{Pr[\,sim matches\,]}}}
    \;\propto\;
    \underbrace{p(\theta \mid b)}_{\substack{\text{kept }\theta\text{'s}\\[2pt]\text{= posterior}}}
  \]}
  \vfill
  \begin{center}
    \small
    $p(b \mid \theta) = \Pr[\,b' = b \mid \theta\,]$:
    the chance one simulation reproduces the data
  \end{center}
\end{frame}

\end{document}
